\chapter{Debugging workflows}

Debugging CI is about making the failure reproducible and visible.
Start by reducing the workflow to the smallest failing case.

\section{Use step logs}
Logs are the primary debugging tool.
Group logs to keep output readable.

\begin{lstlisting}[language=YAML]
- name: Build
  run: |
    echo "::group::Build output"
    ./scripts/build
    echo "::endgroup::"
\end{lstlisting}

\section{Enable debug logging}
GitHub supports two debug flags.
Set them as repository secrets or workflow env vars.

\begin{lstlisting}[language=YAML]
env:
  ACTIONS_STEP_DEBUG: true
  ACTIONS_RUNNER_DEBUG: true
\end{lstlisting}

\section{Rerun jobs}
Rerun failed jobs from the Actions UI.
This saves time when failures are flaky or transient.

\section{Collect diagnostics}
When a failure is intermittent, capture extra artifacts.

\begin{lstlisting}[language=YAML]
- name: Upload logs
  if: ${{ failure() }}
  uses: actions/upload-artifact@v4
  with:
    name: debug-logs
    path: logs/
\end{lstlisting}

\section{Local reproduction}
Reproduce failures locally with the same script the workflow runs.
If you cannot run it locally, move more logic into scripts until you can.

\section{Workflow isolation}
Create a minimal workflow that runs only the failing job.
This reduces noise and isolates the root cause faster.

\section{Manual debug runs}
Add a \code{workflow\_dispatch} entry to run a workflow with specific inputs.
This is useful when you need to debug with a particular branch or parameter set.

\section{Artifacts on failure}
If a job fails intermittently, upload logs or state only on failure.
This keeps successful runs clean while preserving enough detail to investigate.

\begin{checklistbox}{Checklist: debugging}
\begin{itemize}[nosep]
  \item Make the failing step minimal.
  \item Increase logging only when needed.
  \item Capture diagnostics on failure.
  \item Reproduce locally with the same scripts.
\end{itemize}
\end{checklistbox}
